\scanpage{73}
\mspage{H}
\deleted{\unsure{H}}
% archive foliation: 72
c'est un isomorphisme à cause du fait que
$\pi_1(X)\to\pi_1(Y)$ est un isomorphisme
[grâce au th. de pureté sur $Y$]
\margin{$Y$ \uline{régulier}}

Le deuxième homomorphisme s'écrit
\[
H^0(Y,\underline O_Y/\underline O_Y^{p^{\unsure{\nu|1}}})
\longrightarrow
H^0(X,\underline O_X/\underline O_X^{p^{\unsure{\nu|1}}})
\]
[grâce au fait que $X,Y$ sont réduits],
ou aussi, supposant que $k$ est alg. clos donc parfait,
s'interprète comme
\[
f^*:H^0(Y,\underline O_Y^{p^{-\unsure{\nu|1}}}/\underline O_Y)
\longrightarrow
H^0(X,\underline O_X^{p^{-\unsure{\nu|1}}}/\underline O_X)
\]
\margin{$Y$ \unsure{lisse} sur $k$ (!)}

Ce dernier est bijectif grâce à la \uline{normalité} de $X$ et $Y$,
car \deleted{\unsure{si $f$}} avec les notations ci-dessus, on a un
diagramme commutatif
\diagram{\begin{tikzcd}
&T(X)\arrow[dl,"g^*"']\\
T(U)&T(Y)\arrow[l,"\simeq","i^*"']\arrow[u,"f^*"']
\end{tikzcd}}
avec $g^*$ \uline{injective}
[car $\underline O_X^{p^{-\unsure{\nu|1}}}/
\underline O_X^{\unsure{\tau}}$ est \deleted{\unsure{li}}
sans \uline{torsion} \ldots],
d'où $f^*$ et $i^{*-1}g^*$ sont des isom inverses l'un de l'autre.
\margin{$X$ \uline{normal}}
